\documentclass[11pt]{scrartcl}
\usepackage[secthm]{evan}
\usepackage{mathteam}

\title{Demystifying finite fields}
\author{Michael Luo}

\begin{document}

\maketitle

\epigraph{Gauss, this stupid guy, who obviously made a mistake}{Proof by contradiction of Gauss' lemma, Yau Mathcamp}

\tableofcontents
\eject

\section*{Introduction}

Finite fields are a pretty advanced topic, but they come up in short-answer contests surprisingly often. My first exposure to a finite field of non-prime order was \Cref{exa:cubic_discriminant}. The solution was not understandable to me, since I barely knew what a field was. I asked Evan Chen, whose answer was essentially ``read section 60 of Napkin,'' which was not so helpful since section 60 depends on sections 1, 3 through 5, and 53 through 59, for a total of 12 sections.

Later, I attended a five-week course at the \href{https://mathcamp.simis.cn/2024/}{Yau Mathcamp}, finally answering my original question. But you definitely \emph{do not} need five weeks of classes to understand finite fields if you just want to solve contest problems. In fact, I explained \Cref{exa:motivation} to one of my friends in just fifteen minutes today, which inspired this handout.

The goal here is to explain finite fields using as little theory as possible, which comes at the expense of leaving many facts unproved, but allows the reader to learn through examples. In particular, shoot me if I say ``automorphism group.'' We leave only the parts that are relevant to contests, however, we still require experience with modular arithmetic, especially orders, Fermat's little theorem, and modular inverses. If you are unfamiliar with the above, see my orders handout, which may be found on \href{knowingant.github.io/handouts/orders/orders.pdf}{my website}.

\section{What is a field?}

A field is a set where you can add, subtract, multiply, and divide (except by zero), like $\QQ$, $\RR$, and $\CC$. Here, we will work with \emph{finite fields}, which (as their name suggest) have a finite number of elements, such as the integers modulo a prime $p$. The field with $p$ elements is called $\FF_p$ or $\ZZ/p\ZZ$. If you want the precise definition of a field, here it is, but beware: it is boring.

\begin{definition}[Definition for nerds, optional]
    A \emph{field} is a set $F$ along with two binary operations $F^2 \to F$ \emph{addition} and \emph{multiplication}, written $+$ and $\cdot$, respectively, such that the following hold:
    \begin{itemize}
        \ii Addition and multiplication are commutative and associative.
        \ii There exist two elements $0 \neq 1$ such that $a + 0 = a$ and $a \cdot 1 = a$ for all $a \in F$.
        \ii For any $a \in F$, there exists $-a \in F$ such that $a + (-a) = 0$. Furthermore, if $a$ is nonzero, then there exists $a^{-1} \in F$ such that $a \cdot a^{-1} = 1$.
        \ii For any $a,b,c \in F$, $a \cdot (b + c) = a \cdot b + a \cdot c$.
    \end{itemize}
    We often omit the multiplication symbol, writing $ab$ instead of $a \cdot b$.
\end{definition}

See? I told you it's boring.

\section{A motivating example}

\begin{example}[BMT General 2021/21]\label{exa:motivation}
	There exist integers $a$ and $b$ such that $(1 + \sqrt 2)^{12} = a + b\sqrt2$. Compute $ab$ modulo $13$.
\end{example}

\begin{solution}{\Cref{exa:motivation}}
    Let $p = 13$ be a prime. We work in the field
    \[\FF_{p^2} \defeq \{a + b\sqrt2 \mid a,b \in \FF_p\},\]
    so for example, $14 + 34\sqrt 2 = 1 + 8\sqrt2$. This is a field (trust me for now).

    \alert{VERY COMMON MISTAKE!!!} Many people confuse $\FF_{p^2}$ with $\ZZ/p^2\ZZ$. The former is a field (actually ``the'' field, since it's unique, but we'll talk about that later) with $p^2$ elements, the latter is the integers modulo $p^2$. Observe that $\ZZ/p^2\ZZ$ has zero divisors ($p \cdot p = 0$), so it is not a field. Also, in $\FF_{p^2}$, $p = 0$, but in $\ZZ/p^2\ZZ$, $p \neq 0$. Promise yourself you will never make this mistake. Never ever.
    
    Now that we've got that out of the way, let's continue by introducing the supremely cool-sounding Frobenius automorphism. Consider this Frobenius automorphism $\sigma: \FF_{p^2} \to \FF_{p^2}$ given by $\sigma(x) = x^p$. Let me list a bunch of properties of $\sigma$.
    \begin{itemize}
        \ii $\sigma(xy) = \sigma(x)\sigma(y)$: follows since $(xy)^p = x^py^p$.
        \ii $\sigma(x + y) = \sigma(x) + \sigma(y)$: use the binomial theorem and observe that $p \mid \binom pk$ for $1 \le k \le p - 1$.\footnote{This is sometimes known as the ``freshman's dream'' due to the tendency of noobs to use the incorrect identity $(x + y)^n = x^n + y^n$, which is actually true when $n$ is prime and we are working in a field where $n = 0$.}
        \ii The set of fixed points of $\sigma$ is $\FF_p$: if $x \in \FF_p$, then $\sigma(x) - x = x^p - x = 0$ by Fermat's little theorem, but since $x^p - x$ is a polynomial of degree $p$, it has at most $p$ roots in $\FF_{p^2}$, so $\FF_p$ is it.
        \ii Let $\FF_p[x]$ be the set of polynomials with coefficients in $\FF_p$.\footnote{Note that $\FF_p[x]$ is the set of \emph{polynomials}, not the set of polynomial functions. For example, let $P_1(x) = x^2 + x$ and $P_2(x) = 0$. Then $P_1$ and $P_2$ are different polynomials in $\FF_2$, even though $P_1(x) = P_2(x) = 0$ for all $x \in \FF_2$.} For $P \in \FF_p[x]$, $\sigma(P(x)) = P(\sigma(x))$: write $P(x) = a_nx^n + a_{n - 1}x^{n - 1} + \dots + a_0x^0$, then
        \begin{align*}
            \sigma(P(x)) &= \sigma(a_nx^n + a_{n - 1}x^{n - 1} + \dots + a_0x^0) \\
            &= \sigma(a_nx^n) + \sigma(a_{n - 1}x^{n - 1}) + \dots + \sigma(a_0x^0) \\
            &= \sigma(a_n)\sigma(x)^n + \sigma(a_{n - 1})\sigma(x)^{n - 1} + \dots + \sigma(a_0)\sigma(x)^0 \\
            &= a_n\sigma(x)^n + a_{n - 1}\sigma(x)^{n - 1} + \dots + a_0\sigma(x)^0 \\
            &= P(\sigma(x)).
        \end{align*}
        \ii If $P \in \FF_p[x]$ has $\alpha$ as a root, it also has $\sigma(\alpha)$ as a root: $P(\sigma(\alpha)) = \sigma(P(\alpha)) = \sigma(0) = 0$.
    \end{itemize}
    
    Now
    \[(1 + \sqrt2)^{12} = \frac{\sigma(1 + \sqrt2)}{1 + \sqrt2} = \sigma(1 + \sqrt2)(\sqrt2 - 1),\]
    so we just need to compute $\sigma(1 + \sqrt2)$. Now since $1 + \sqrt2$ is a root of the polynomial
    \[(x - (1 + \sqrt2))(x - (1 - \sqrt2)) = x^2 - 2x - 1,\]
    $\sigma(1 + \sqrt2)$ is either $1 + \sqrt2$ or $1 - \sqrt2$. But since $1 + \sqrt2 \not\in \FF_p$, it's not a fixed point of $\sigma$, so $\sigma(1 + \sqrt2) = 1 - \sqrt2$. So
    \[\sigma(1 + \sqrt2)(\sqrt2 - 1) = (1 - \sqrt2)(\sqrt2 - 1) = 10 + 2\sqrt2,\]
    which gives the answer $7$.
\end{solution}

\begin{comment}
\begin{remark*}
    Note that we \emph{never used the condition} that $\FF_{p^2}$ was the set of $a + b\sqrt2$, actually, we could have made $\FF_{p^2}$ be the set of $6a + 9b\sqrt2$, and it wouldn't have changed the solution at all!
\end{remark*}
\end{comment}

\section{What \texorpdfstring{$\FF_{p^n}$}{} really means}

Sorry, this small amount of theory is actually necessary. I swear, it's not a lot.

Let $p = 13$ again, as in \Cref{exa:motivation}. Remember I told you
\[\FF_{p^2} = \{a + b\sqrt2 \mid a,b \in \FF_p\}\]
is a field? Let me actually explain why.

First, I claim the set described above can be thought of as
\[\FF_p[x]/(x^2 - 2).\]
Wait, what does this notation mean? You know how $\ZZ/5$ is the integers modulo $5$? In the same way, $\FF_p[x]/(x^2 - 2)$ is the set of polynomials with coefficients in $\FF_p$, ``modulo $x^2 - 2$.'' So basically, we can add and subtract multiples of $x^2 - 2$ at will. For example, in $\FF_p[x](x^2 - 2)$, we have
\[x^3 = x^3 - x(x^2 - 2) = 2x.\]
Why is this the same concept as above? Well, just look at these two analogous calculations: $x$ is basically $\sqrt2$.
\begin{multline*}
    (3 + 5\sqrt2)(1 + 8\sqrt2) = 3 + 29\sqrt2 + 40\sqrt2^2 = 3 + 3\sqrt2 + \sqrt2^2 = 3 + 3\sqrt2 + 2 = 5 + 3\sqrt2 \\
    (3 + 5x)(1 + 8x) = 3 + 29x + 40x^2 = 3 + 3x + x^2 = 3 + 3x + 2 = 5 + 3x.
\end{multline*}

I'll now convince you this is a field. Note that $\ZZ/5$ is a field, since $5$ is prime. Similarly, $x^2 - 2$ is ``prime'' in $\FF_p[x]$ (it is a irreducible polynomial); since $2$ is a quadratic nonresidue modulo $13$, we see $x^2 - 2$ has no roots. Therefore, $\FF_p[x]/(x^2 - 2)$ is a field, just like $\ZZ/5$ is a field. But $\FF_p[x]/(x^2 - 1)$ is not a field, just like $\ZZ/6$ is not a field. (This is since $(x - 1)(x + 1) = x^2 - 1 = 0$.) But you can decompose $\ZZ/6 = \ZZ/2 \times \ZZ/3$, so similarly, you can decompose
\[\FF_p[x]/(x^2 - 1) = \FF_p[x]/(x - 1) \times \FF_p[x]/(x + 1) = \FF_p \times \FF_p = \FF_p^2\]
by the ``Chinese remainder theorem.''

\begin{remark*}
    The reason we usually write $\FF_{p^2}$ instead of $\FF_p[x]/(x^2 - 2)$ is that \emph{the exact choice of the polynomial $x^2 - 2$ does not matter}! If $P$ and $Q$ are irreducible in $\FF_p[x]$ with degree $2$, then the fields
    \[\FF_p[x]/P(x) \qquad \text{and} \qquad \FF_p[x]/Q(x)\]
    are isomorphic. More generally, \emph{all finite fields with the same number of elements are isomorphic!}
    
    For every prime power $q$, there is exactly one (up to isomorphism) finite field with $q$ elements. Proof sketch: all finite fields must have prime characteristic, so they must have prime power cardinality. All fields with cardinality $q = p^a$ are isomorphic to $\{x \in \ol{\FF_p} \mid x^q - x = 0\}$.
\end{remark*}

So the whole point of $\FF_{p^n}$ is basically that you can deal with a degree-$n$ irreducible polynomial with coefficients in $\FF_p$.

Also, note that $\FF_{p^n}$ satisfies a form of Fermat's little theorem: $x^{p^n} = x$ for all $x \in \FF_{p^n}$. Also, a nice fact is that there are primitive roots in $\FF_{p^n}$: that is, there are elements of order $p^n - 1$. The proofs for these facts are exactly the same as the analogous facts in $\FF_p$. The point I'm trying to drive home here is that finite fields are really, \emph{really} nice.

We continue with examples!

\section{A bunch of examples}

Try the examples below before reading the solutions.

\begin{example}[BMT Discrete 2023/10]\label{exa:bprp}
	Let $\alpha$ be the positive real root of $x^2 - 3x - 2$. Compute $\floor{\alpha^{1000}}$ modulo $997$.
\end{example}

\begin{solution}{\Cref{exa:bprp}}
    Let $p = 997$. Let $\beta$ be the other root of $x^2 - 3x - 2$. From the quadratic formula,
    \[\alpha = \frac{3 + \sqrt{17}}{2} \qquad \beta = \frac{3 - \sqrt{17}}{2}.\]
    Since $\abs{\beta} < 1$, $0 < \beta^{1000} < 1$, so we just want
    \[\alpha^{1000} + \beta^{1000} - 1.\]
    We work in the field
    \[\FF_{p^2} = \{a + b\sqrt{17} \mid a,b \in \FF_p\}.\]
    Let $\sigma$ be the Frobenius automorphism sending $x$ to $x^p$. Then since $\alpha$ and $\beta$ are roots of the polynomial $x^2 - 3x - 2$, and $\alpha,\beta \not\in \FF_p$, $\sigma(\alpha) = \beta$ and $\sigma(\beta) = \alpha$. Therefore
    \[\alpha^{1000} + \beta^{1000} - 1 = \alpha^3\sigma(\alpha) + \beta^3\sigma(\beta) - 1 = \alpha\beta(\alpha^2 + \beta^2) - 1 = -27 = 970. \qedhere\]
\end{solution}

\begin{example}[Hong Kong 2019/1/4]\label{exa:china}
    How many primes $p < 100$ satisfy $p \mid \floor{(2 + \sqrt 5)^p} - 2^{p + 1}$?
\end{example}

\begin{solution}{\Cref{exa:china}}
    If $p = 2$, we fail immediately. Otherwise, $p$ is odd, so we want
    \[(2 + \sqrt5)^p + (2 - \sqrt5)^p - 2^{p + 1}.\]
    But since $\sigma(a) + \sigma(b) = \sigma(a + b)$, this is
    \[4^p - 2^{p + 1} = 2^{p + 1}(2^{p - 1} - 1) \equiv 0 \pmod p,\]
    so all odd $p$ below $100$ work. So the answer is $24$.
\end{solution}

\begin{example}[HMMT Algebra and Number Theory 2017/9]\label{exa:fibonacci}
    Compute the period of the Fibonacci sequence modulo $127$.
\end{example}

\begin{solution}{\Cref{exa:fibonacci}}
    Let $p = 127$, $\varphi = \frac{1 + \sqrt5}2$, and $\psi = \frac{1 - \sqrt5}2$. We apply Binet's formula
    \[F_n = \frac{\varphi^n - \psi^{n}}{\sqrt5}.\]
    We work in the field $\FF_{p^2}$. Since $(\frac{5}{127}) = -1$, $\varphi,\psi \not\in \FF_p$. So $\varphi$ and $\psi$ are roots of the polynomial $x^2 - x - 1$, $\sigma(\varphi) = \psi$ and $\sigma(\psi) = \varphi$. We compute
    \begin{align*}
        F_p &= \frac{\varphi^p - \psi^p}{\sqrt5} = \frac{\psi - \varphi}{\sqrt5} = -1 \\
        F_{p + 1} &= \frac{\varphi\psi - \varphi\psi}{\sqrt5} = 0 \\
        F_{2p + 1} &= \frac{\varphi\psi^2 - \varphi^2\psi}{\sqrt5} = 1 \\
        F_{2p + 2} &= \frac{\varphi^2\psi^2 - \varphi^2\psi^2}{\sqrt5} = 0,
    \end{align*}
    so $F_{2p + 2} = F_{p + 1} + F_p = -1$ and $F_{2p + 3} = F_{2p + 2} + F_{2p + 1} = 1$. Therefore
    \[(F_0,F_1) = (F_{256},F_{257}) \neq (F_{128},F_{129}),\]
    so the period divides $256$ but not $128$, so it is exactly $256$.
\end{solution}

\begin{example}[HMMT Guts 2023/31]\label{exa:cubic_discriminant}
    Note that $2017$ is prime. Compute
    \[\prod_{i = 0}^{2016}(i^3 - i - 1)^2 \pmod{2017}\]
    given it's not zero.
\end{example}

\begin{solution}{\Cref{exa:cubic_discriminant}}
    Let $p = 2017$ and $P(x) = x^3 - x - 1$. Since the answer is nonzero, $P$ has no roots modulo $p$, so it is irreducible modulo $p$. Therefore, let
    \[x^3 - x - 1 = (x - a)(x - b)(x - c) \qquad \text{where $a,b,c \in \FF_{p^3}$.}\]
    Also, since everything is squared, signs don't matter. Since $(x - 0)(x - 1)\dots(x - (p - 1)) = x^p - x$, we want
    \[\biggl(\prod_{i \in \FF_p}(i - a)(i - b)(i - c)\biggr)^2 = \bigl((a^p - a)(b^p - b)(c^p - c)\bigr)^2.\]
    Since $a,b,c \not\in \FF_p$ and they are all roots of $P$, $(a^p,b^p,c^p)$ is a nontrivial cycle of $(a,b,c)$, so this becomes
    \[\bigl((a - b)(b - c)(c - a)\bigr)^2\]

    This is the discriminant of $x^3 - x - 1$, which is $-4 \cdot (-1)^3 - 27(-1)^2 = -23 \equiv 1994 \pmod p$. Alternatively, we may take the derivative: a double application of product rule gives
    \[P'(x) = 3x^2 - 1 = (x - a)(x - b) + (x - b)(x - c) + (x - c)(x - a),\]
    so
    \[P'(a) = (a - b)(a - c) \qquad P'(b) = (b - c)(c - a) \qquad P'(c) = (c - a)(c - b).\]
    Therefore
    \begin{align*}
        \bigl((a - b)(b - c)(c - a)\bigr)^2 &= -P'(a)P'(b)P'(c) \\
        &= (1 - 3a^2)(1 - 3b^2)(1 - 3c^2) \\
        &= -27P\left(\frac1{\sqrt 3}\right)P\left(-\frac1{\sqrt 3}\right) \\
        &= -27\left(1 - \frac2{3\sqrt 3}\right)\left(1 + \frac2{3\sqrt 3}\right) \\
        &= -23,
    \end{align*}
    so the answer is $1994$.
\end{solution}

\begin{example}[Brazil 2017/6]\label{exa:brazil}
    If $p \mid a^3 - 3a + 1$ and $p \neq 3$, show that $p \equiv \pm 1 \pmod 9$.
\end{example}

\begin{solution}{\Cref{exa:brazil}}
    This one is rather stupid. Work modulo $p$, then let $u \in \FF_{p^2}$ be a root of the polynomial $x^2 - ax + 1$, so $u + \frac1u = a$. Then
    \[a^3 - 3a + 1 = \left(u + \frac1u\right)^3 - 3\left(u + \frac1u\right) + 1 = u^3 + \frac1{u^3} + 1 = 0,\]
    so $u^6 + u^3 + 1 = 0$, so $u^9 = 1$ but $u^3 \neq 1$. Therefore, $u$ is an element of order $9$ in $\FF_{p^2}^\times$, so $9 \mid p^2 - 1$, or $p \equiv \pm 1 \pmod 9$.
\end{solution}

How does the above solution break if $p = 3$?

\section{Exercises}

\begin{exercise}[HMMT Algebra and Number Theory 2015/8]
    Find the number of ordered pairs of integers $(a,b) \in \{1,2,\dots,35\}^2$ (not necessarily distinct) such that there exist polynomials $f$, $P$, and $Q$ with integer coefficients such that
    \[f(x)^2 - (ax + b) = (x^2 + 1)P(x) + 35Q(x).\]
\end{exercise}

\begin{exercise}[BMT Discrete 2022/8]
    Define the two sequences $a_0,a_1,a_2,\dots$ and $b_0,b_1,b_2,\dots$ by $a_0 = 3$ and $b_0 = 1$ with the recurrence relations $a_{n + 1} = 3a_n + b_n$ and $b_{n + 1} = 3b_n - a_n$ for all nonnegative integers $n$. Compute $a_{32}$ and $b_{32}$ modulo $31$.
\end{exercise}

\begin{exercise}[BMT Individual 2017/20]
    Compute
    \[\sum_{k = 0}^{15}(-1)^k\left(2\cos\left( \frac{k\pi}{16} \right)\right)^{560} \pmod{17}.\]
\end{exercise}

\begin{exercise}[HMMT Algebra and Number Theory 2020/10]
    Let $p = 101$. We say $f \in \FF_p[x]$ is \emph{lucky} if it has degree at most $1000$ and there exist $g,h \in \FF_p[x]$ such that
    \[f(x) = g(x)(x^{1001} - 1) + h(x)^p - h(x)\]
    in $\FF_{101}[x]$. Compute the number of lucky polynomials.
\end{exercise}

\end{document}